\section{Max-value Entropy Search}
%
Entropy search methods use an information-theoretic perspective to select where to evaluate. They
%
 find a query point that maximizes the information about the location $\vx_* = \argmax_{\vx \in \mathfrak X} f(x)$ whose value $y_* = f(\vx_*)$ achieves the global maximum of the function $f$. Using the negative differential entropy of $p(\vx_* | D_t)$ to characterize the uncertainty about $\vx_*$, ES and PES use the acquisition functions
 %
 \begin{align}
   %
&\alpha_t(x) = 
I(\{\vx,y\}; \vx_*\mid D_t) \\
&\;\;\;\;= H\left(p(\vx_*\mid D_t)\right)- \Ex\left[H(p(\vx_*\mid D_t\cup \{\vx, y\})) \right] \label{eq:es}\\
%
&\;\;\;\; = H(p(y\mid D_t,\vx))- \Ex\left[H(p(y\mid D_t, \vx, \vx_*)) \right] \label{eq:pes}.
%
\end{align}
%
ES uses formulation~\eqref{eq:es}, in which the expectation is over $p(y| D_t, \vx)$, while PES uses the equivalent, symmetric formulation~\eqref{eq:pes}, where the expectation is over $p(\vx_* | D_t)$. Unfortunately, both $p(\vx_* | D_t)$ and its entropy is analytically intractable and have to be approximated via expensive computations.
Moreover, the optimum may not be unique, adding further complexity to this distribution.
%
%

We follow the same information-theoretic idea but propose a much cheaper and more robust objective to compute. Instead of measuring the information about the argmax $\vx_*$, we use the information about the \emph{maximum value} $y_* = f(\vx_*)$. 
Our acquisition function is the gain in mutual information between the maximum $y_*$ and the next point we query, which can be approximated analytically by evaluating the entropy of the predictive distribution:
\begin{align}
&\alpha_t(x) = 
I(\{\vx,y\}; y_*\mid D_t) \\
&= H(p(y\mid D_t, \vx)) - \Ex[H(p(y\mid D_t, \vx, y_*)) ] \label{eq:mes} \\
& \approx \frac{1}{K}\sum_{y_*\in Y_*} \left[\frac{\gamma_{y_*}( \vx)\psi(\gamma_{y_*}( \vx))}{2\Psi(\gamma_{y_*}(\vx))} - \log(\Psi(\gamma_{y_*}( \vx))) \right] \label{eq:mesapprox}
\end{align}
%
where $\psi$ is the probability density function and $\Psi$ the cumulative density function of a normal distribution, and $\gamma_{y_*}(\vx) = \frac{y_* - \mu_t(\vx)}{\sigma_t(\vx)}$. %
The expectation in Eq.~\eqref{eq:mes} is over $p(y_*|D_n)$, which is approximated using Monte Carlo estimation by sampling a set of $K$ function maxima. %
 Notice that the probability in the first term $p(y | D_t, \vx)$ is a Gaussian distribution with mean $\mu_t(\vx)$ and variance $k_t(\vx,\vx)$. The probability in the second term $p(y | D_n, \vx, y_*)$ is a truncated Gaussian distribution: given $y_*$, the distribution of $y$ needs to satisfy $y < y_*$.  Importantly, while ES and PES rely on the expensive, $d$-dimensional distribution $p(\vx_* | D_t)$, here, we use the one-dimensional $p(y_*|D_n)$, which is computationally much easier.

It may not be immediately intuitive that the \emph{value} should bear sufficient information for a good search strategy. Yet, the empirical results in Section~\ref{sec:exp} will demonstrate that this strategy is typically at least as good as ES/PES. From a formal perspective,  \citet{wang2016est} showed how an estimate of the maximum value implies a good search strategy (EST). Indeed, Lemma~\ref{lem:equivalence} will make the relation between EST and a simpler, degenerate version of MES explicit.

Hence, it remains to determine how to sample $y_*$. We propose two strategies: (1) sampling from an approximation via a Gumbel distribution; and (2) sampling functions from the posterior Gaussian distribution and maximizing the functions to obtain samples of $y_*$. We present the MES algorithm in Alg.~\ref{alg:mes}.%

\begin{algorithm} %
  \caption{Max-value Entropy Search (MES)%
  }\label{alg:mes}
  \begin{algorithmic}[1]
  %
  %
  %
    %
    %
    \FUNCTION{MES\,($f, D_0$)}
      \FOR{$t = 1,\cdots, T $}
      %
      \STATE $\alpha_{t-1}(\cdot)\gets$\textsc{Approx-MI\,}($D_{t-1}$)%
      \STATE $\vx_t\gets \argmax_{\vx\in\mathfrak X}{\alpha_{t-1}(\vx)}$ 
      \STATE $y_t\gets f(\vx_t) + \epsilon_t, \epsilon_t\sim\mathcal N(0,\sigma^2)$
      \STATE $\mathfrak D_t \gets D_{t-1}\cup \{\vx_t,y_t\}$
      %
      \ENDFOR
      \ENDFUNCTION
      \item[]
    \FUNCTION{Approx-MI\,($D_t$)}%
    \IF{Sample with Gumbel}
    \STATE approximate $\Pr[\hat{y}_* < y]$ with $\mathcal G(a,b)$
    \STATE sample a $K$-length vector $\vr\sim \text{Unif}([0,1])$
    \STATE $\vy_* \gets a-b\log(-\log \vr)$
    \ELSE
    \FOR{$i=1,\cdots,K$}
    \STATE sample $\tilde f \sim GP(\mu_t,k_t \mid D_t)$
    \STATE $y_{*(i)} \gets \max_{\vx\in\mathfrak X}{\tilde f(\vx)}$
    \ENDFOR 
    \STATE $\vy_* \gets [y_{*(i)}]_{i=1}^K$
    \ENDIF
    \STATE \textbf{return} $\alpha_t(\cdot)$ in Eq.~\eqref{eq:mesapprox}
    \ENDFUNCTION
  \end{algorithmic}
\end{algorithm}

\subsection{Gumbel Sampling} %
\label{spec:gumbel}
The marginal distribution of $f(x)$ for any $x$ is a one-dimensional Gaussian, and hence the distribution of $y^*$ may be viewed as the maximum of an infinite collection of dependent Gaussian random variables.
%
Since this distribution is difficult to compute, we make two simplifications. First, we replace the continuous set $\mathfrak X$ by a discrete (finite), dense subset $\hat{\mathfrak X}$ of representative points. If we select $\hat{\mathfrak X}$ to be an $\epsilon$-cover of $\mathfrak X$ and the function $f$ is Lipschitz continuous with constant $L$, then we obtain a valid upper bound on $f(\mathfrak X)$ by adding $\epsilon L$ to any upper bound on $f(\hat{\mathfrak X})$.

Second, we use a ``mean field'' approximation and treat the function values at the points in $\hat{\mathfrak X}$ as independent. This approximation tends to over-estimate the maximum; this follows from Slepian's lemma if $k(x,x') \geq 0$. Such upper bounds still lead to optimization strategies with vanishing regret, whereas lower bounds may not \citep{wang2016est}.

%

We sample from the approximation $\hat{p}(y^*|D_n)$ via its cumulative distribution function (CDF) $\widehat{\Pr}[y_* < z] = \prod_{\vx\in\hat{\mathfrak X}} \Psi(\gamma_z(\vx))$. That means we sample $r$ uniformly from $[0,1]$ and find $z$ such that $\Pr[y_* < z] = r$. A binary search for $z$ to accuracy $\delta$ requires $O(\log\tfrac1\delta)$ queries to the CDF, and each query takes $O(|\hat{\mathfrak X}|)\approx O(n^d)$ time, so we obtain an overall time of $O(M|\hat{\mathfrak X}| \log \frac{1}{\delta})$ for drawing $M$ samples.

\iffalse
\sj{ clarify: Let $\tilde y_*$ be an upper bound on the function values of $\hat{\mathfrak X}$. We have $\Pr[\hat{y}_* < y] = \prod_{\vx\in\hat{\mathfrak X}} \Psi(\gamma_y(\vx)))$. With the cumulative probability, we can sample $\hat{y}_*$ by sampling $r$ uniformly from $[0,1]$ and find $y$ such that $\Pr[\hat{y}_* < y] = r$.}
We can find such a $y$ by binary search. Each query to $\Pr[\hat{y}_* < y]$ takes time $O(|\hat{\mathfrak X})|\approx O(n^d)$, where $d$ is the input dimension, so we obtain an overall time of $O(|\hat{\mathfrak X}| \log \frac{1}{\delta})$ for drawing $M$ samples with accuracy $\delta$ of the binary search.
\fi
%

\iffalse
In \cite{wang2016est}, the authors estimated an upper bound on the function $f$ by assuming independence among the Gaussian values of a discrete set of input candidates. Here we adopt a similar strategy by selecting a representative finite set of input candidates $\hat{\mathfrak X}$ from the continuous space $\mathfrak X$. Theoretically if we select $\hat{\mathfrak X}$ to be the $\epsilon$-covering of $\mathfrak X$ and the function $f$ is Lipschitz continuous, we can obtain an upper bound on $f(\mathfrak X)$ by adding $\epsilon L$ to the upper bound on $f(\hat{\mathfrak X})$.


Let $\tilde y_*$ be an upper bound on the function values of $\hat{\mathfrak X}$. We have $\Pr[\hat{y}_* < y] = \prod_{\vx\in\hat{\mathfrak X}} \Psi(\gamma_y(\vx)))$. With the cumulative probability, we can sample $\hat{y}_*$ by sampling $r$ uniformly from $[0,1]$ and find $y$ such that $\Pr[\hat{y}_* < y] = r$. However, each query of $\Pr[\hat{y}_* < y]$ takes $O(|\hat{\mathfrak X})|\approx O(n^d_x)$ where $d_x$ is the dimension of the input, and numerically solving $\Pr[\hat{y}_* < y] = r$ takes $O(|\hat{\mathfrak X}| \log \frac{1}{\delta})$ with a binary search strategy, where $\delta$ is the accuracy for the solution $y$. So a naive way of sampling $M$ number of $\hat y_*$ inevitably takes $O(M|\hat{\mathfrak X}|\log \frac{1}{\delta})$. 
\fi


To sample more efficiently, we propose a $O(M+|\hat{\mathfrak X}|\log \frac{1}{\delta})$-time strategy, by approximating the CDF by a Gumbel distribution: $\widehat{\Pr}[y_* < z]\approx \mathcal G(a,b) = e^{-e^{-\frac{z-a}{b}}}$. This choice is motivated by the Fisher-Tippett-Gnedenko theorem~\cite{fisher1930genetical}, which states that the maximum of a set of i.i.d. Gaussian variables is asymptotically described by a Gumbel distribution (see the appendix for further details). This does not in general extend to non-i.i.d. Gaussian variables, but we nevertheless observe that in practice, this approach yields a good and fast approximation.

We sample from the Gumbel distribution via the Gumbel quantile function: we sample $r$ uniformly from $[0,1]$, and let the sample be $y = \mathcal G^{-1}(a,b) = a-b\log(-\log r)$. We set the appropriate Gumbel distribution parameters $a$ and $b$ by percentile matching and solve the two-variable linear equations $a-b\log(-\log r_1) = y_1$ and $a-b\log(-\log r_2) = y_2$, where $\Pr[{y}_* < y_1]=r_1$ and $\Pr[{y}_* < y_2]=r_2$. In practice, we use $r_1=0.25$ and $r_2=0.75$ so that the scale of the approximated Gumbel distribution
is proportional to the interquartile range of the CDF $\hat{\Pr}[y_* < z]$. %


\iffalse
To sample more efficiently, we propose a $O(M+|\hat{\mathfrak X}|\log \frac{1}{\delta})$ strategy by approximating $\Pr[\hat{y}_* < y]$ with a Gumbel distribution $\Pr[\hat{y}_* < y]\approx \mathcal G(a,b) = e^{-e^{-\frac{y-a}{b}}}$. To generate samples from a Gumbel distribution, we use the Gumbel quantile function and set the sample to be $y = \mathcal G^{-1}(a,b) = a-b\log(-\log r)$ where $r$ is uniformly sampled from the interval $[0,1]$. To find the appropriate Gumbel distribution parameters $a$ and $b$, we use percentile matching and solve the two-variable linear equations $a-b\log(-\log r_1) = y_1$ and $a-b\log(-\log r_2) = y_2$, where $\Pr[\hat{y}_* < y_1]=r_1$ and $\Pr[\hat{y}_* < y_2]=r_2$. In practice we set $r_1=0.25$ and $r_2=0.75$ so that the interquartile range will determine the variance of the approximated Gumbel distribution.

Notice that for the Gaussian process prior, using Gumbel distribution for the approximation of the maximum of i.i.d Gaussians is asymptotically accurate by the Fisher-Tippett-Gnedenko theorem~\cite{fisher1930genetical}. The Fisher-Tippett-Gnedenko theorem in general does not hold for independent and differently distributed Gaussians, but in practice we still find it useful in approximating $\Pr[\hat{y}_* < y]$\footnote{See the appendix for more details.}. 
\fi

\subsection{Sampling $y_*$ via Posterior Functions}
For an alternative sampling strategy we follow~\cite{hernandez2014predictive}: we draw functions from the posterior GP and then maximize each of the sampled functions.
%
Given the observations $D_t = \{(\vx_\tau, y_\tau)_{\tau=1}^t\}$, we can approximate the posterior Gaussian process using a 1-hidden-layer neural network $\tilde f(\vx) = \va_t\T \vphi(\vx)$ where $\vphi(x)\in \R^D$ is a vector of feature functions~\cite{neal96,rahimi2007random}
%
and the Gaussian weight $\va_t\in\R^D$ is distributed according to a multivariate Gaussian $\mathcal N(\vnu_t, \mSigma_t)$.

\textbf{Computing $\vphi(x)$.} By Bochner's theorem~\cite{rudin2011fourier}, the Fourier transform $\hat k$ of a continuous and translation-invariant kernel $k$ is guaranteed to be a probability distribution. Hence we can write the kernel of the GP to be $k(\vx,\vx') = \Ex_{\ww\sim\hat k(\ww)}[e^{i\ww\T (\vx-\vx')}] =\Ex_{c\sim U[0,2\pi]} \Ex_{\hat k} [2 \cos(\ww\T \vx + c)\cos(\ww\T \vx' + c)]$ and approximate the expectation by $k(\vx,\vx')\approx \vpp\T(\vx)\vpp(\vx')$ where $\phi_i(\vx) = \sqrt{\frac2D} \cos(\ww_i\T \vx + c_i)$, $\ww_i\sim \hat \kk(\ww)$, and $c_i \sim U[0,2\pi]$ for $i = 1,\dots, D$.

\textbf{Computing $\vnu_t, \mSigma_t$.} By writing the GP as a random linear combination of feature functions $\va^T_t\vphi(\vx)$, we are defining the mean and covariance of the GP to be $\mu_t(\vx) = \vnu\T \vphi(\vx)$ and $k(\vx,\vx') = \vphi(\vx)\T\mSigma_t\vphi(\vx')$. Let $Z = [z_1, \cdots, z_t] \in \R^{D\times t}$, where $z_\tau \coloneqq \vpp (\vx_\tau) \in \R^D$. 
The GP posterior mean and covariance in Section~\ref{ssec:gp} become
$\mu_t(\vx) = z\T Z (Z\T Z + \sigma^2 \mI)^{-1} \vy_t$ and 
$k_t (\vx, \vx')= z\T z' - z\T Z (Z\T Z + \sigma^2 \mI)^{-1} Z\T z'$. 
Because $Z(Z\T Z + \sigma^2 \mI)^{-1} = (ZZ\T +\sigma^2 \mI)^{-1}Z$, we can simplify the above equations and obtain $\vnu_t = \sigma^{-2}\mSigma_t Z_t\vy_t $ and $\mSigma_t = (Z Z\T \sigma^{-2} +\mI)^{-1}$.
%


To sample a function from this random 1-hidden-layer neural network, we sample $\tilde \va$ from $\mathcal N(\vnu_t, \mSigma_t)$ and construct the sampled function $\tilde f = \tilde \va\T \vphi(\vx)$. Then we optimize $\tilde f$ with respect to its input to get a sample of the maximum of the function $\max_{\vx\in\mathfrak X} \tilde f(\vx)$. 



\hide{
We first review Mercer's theorem~\cite{mercer1909functions}, which builds the foundation of kernel methods.
\begin{thm}[Mercer~\cite{mercer1909functions}]
Any kernel $\kk:\mathfrak X \times \mathfrak X \rightarrow \R$ satisfying $\int \kk(x,x')f(x)f(x')\dif x \dif x' \geq 0$ for all $L_2(\mathfrak X)$ measurable functions $f$ can be expanded into 
\begin{align}
\kk(x,x') = \sum_j \lambda^{(j)} \phi_j(x)\phi_j(x'), 
\end{align}
 where $\lambda_j$ and $\phi_j$ are orthonormal on $L_2(\mathfrak X)$.
 

\end{thm}
Without loss of generality, we assume $\sum_j \lambda^{(j)} = 1, p(\lambda^{(j)}) = \lambda^{(j)}$. We can get a low variance approximation of $\kk(x,x')$ by sampling $\lambda\sim p(\lambda)$ and the corresponding random feature $\phi_\lambda$.
\begin{align}
\kk(x,x') = \Ex_{\lambda}[\phi_\lambda(x)\phi_{\lambda}(x')] \approx \frac{1}{D} \sum_{i=1}^D \phi_{\lambda_i}(x)\phi_{\lambda_i}(x') = \frac{1}{D}\vpp(x)\T\vpp(x'), 
\end{align}

where $\vpp(x)\in \R^D$ is the random feature vector for $x$. Let $z \coloneqq \frac{1}{\sqrt{D}} \vpp (x), z_\tau \coloneqq \frac{1}{\sqrt{D}} \vpp (x_\tau) \in \R^D$ and $Z = [z_1, \cdots, z_t] \in \R^{D\times t}$.

The GP posterior mean and covariance become
\begin{align}
\mu_t(x) &= z\T Z (Z\T Z + \sigma^2 \mI)^{-1} \vy_t, \label{eq:gpprimal1}\\
\kk_t (x)&= z\T z - z\T Z (Z\T Z + \sigma^2 \mI)^{-1} Z\T z. \label{eq:gpprimal2}
\end{align}

%
Because $Z(Z\T Z + \sigma^2 \mI)^{-1} = (ZZ\T +\sigma^2 \mI)^{-1}Z$, Eq.~\eqref{eq:gpprimal1} and Eq.~\eqref{eq:gpprimal2} can be simplified to

%
\begin{align}
\Sigma_t & = (Z Z\T \sigma^{-2} +\mI)^{-1}\\
\mu_t(x) &= \sigma^{-2}z\T \Sigma_t Z \vy_t  \\
\kk_t(x) &= z\T \Sigma_t z %
\end{align}
%
}


\subsection{Relation to Other BO Methods}
As a side effect, our new acquisition function draws connections between ES/PES and other popular BO methods. The connection between MES and ES/PES follows from the information-theoretic viewpoint; the following lemma makes the connections to other methods explicit.
%
%
\begin{lem}\label{lem:equivalence}
  The following methods are equivalent:
  \vspace{-5pt}
\begin{enumerate}\setlength{\itemsep}{-1pt}
\item MES, where we only use a single sample $y_*$ for $\alpha_t(x)$;
\item EST with $m = y_*$;
\item GP-UCB with $\beta^{\frac12} = \min_{\vx\in\mathfrak X} \frac{y_*-\mu_{t} (\vx)}{\sigma_{t} (\vx)}$;
\item PI with $\theta=y_*$.
\end{enumerate}
\end{lem}
This equivalence no longer holds if we use $M>1$ samples of $y_*$ in MES.
\begin{proof}
  The equivalence among 2,3,4 is stated in Lemma 2.1 in~\cite{wang2016est}. What remains to be shown is the equivalence between 1 and 2. When using a single $y_*$ in MES,
  the next point to evaluate is chosen by maximizing $\alpha_t(\vx) =\gamma_{y_*}( \vx)\frac{\psi(\gamma_{y_*}( \vx))}{2\Psi(\gamma_{y_*}(\vx))} - \log(\Psi(\gamma_{y_*}( \vx)))$ and $\gamma_{y_*} = \frac{y_* - \mu_t(\vx)}{\sigma_t(\vx)}$. For EST with $m=y_*$, the next point to evaluate is chosen by minimizing $\gamma_{y_*}( \vx)$. Let us define a function $g(u) = u\frac{\psi(u)}{2\Psi(u)} - \log(\Psi(u))$. Clearly, $\alpha_t(\vx) = g(\gamma_{y_*}(\vx))$. Because $g(u)$ is a monotonically decreasing function, maximizing $g(\gamma_{y_*}(\vx))$ is equivalent to minimizing $\gamma_{y_*}(\vx)$. Hence 1 and 2 are equivalent.
\end{proof}

%

%


